%#########################################
% For 140601 GitHub Kaigi
% title: GitHubで行うreproducible research
%#########################################
\documentclass[dvipdfmx, 14pt]{beamer}
\usepackage{pxjahyper, otf, mediabb, tikz, wrapfig}
\usepackage{pifont}
\usetheme{default}
\renewcommand{\kanjifamilydefault}{\gtdefault}
\renewcommand{\familydefault}{\sfdefault} % 欧文書体をHelveticaに
\setbeamertemplate{navigation symbols}{}

\setbeamertemplate{background canvas}{\includegraphics[height=\paperheight]{bg}}


\setbeamerfont{title}{size=\huge, series=\bf}
\setbeamercolor{title}{fg=black}
\setbeamercolor{normal text}{fg=black, bg=custom.bgc}
\setbeamercolor{frametitle}{fg=black}

\definecolor{custom.fc}{HTML}{F4F6F7}
\definecolor{custom.bgc}{HTML}{4798DD}

\graphicspath{{../../src/}}

%#######*****####### titile #######*****#######***
\title{GitHubで行う\\reproducible research}
\author[瓜生]{\href{http://urush.postach.io}{瓜生真也}}
\institute[横国大院・環情]{横浜国立大学大学院 環境情報学府}
\subject{LT slide for GitHubKaigi@Shibuya}
\date{\scriptsize GitHub Kaigi(@Shibuya), June 1, 2014\\
\vspace{1em}
\href{http://github.com/uribo/rep-res-guideline}{github.com/uribo}
}
%#######*****#######*****#######*****#######***
\begin{document}

\frame{\titlepage}

\frame{
  \begin{center}
  \Huge 「論文」と聞いて\\どういう印象を\\持ちますか？
  \end{center}}

{\usebackgroundtemplate{\includegraphics[width=\paperwidth]{gg-suggest}}%
\frame{}}%なぜ論文に再現性が必要なのか

{\usebackgroundtemplate{\includegraphics[width=\paperwidth]{gg-trends}}%
\frame{}}%尤ももっともらしい成果を示すために

\frame{
\frametitle{Reproducible Research (RR)}
  {\Large 再現可能な研究}\\
  \vspace{0.5em}
  論文に書かれている方法、条件で行えば、\\同様の結果が得られる
  \begin{wrapfigure}{o}[60pt]{60mm}
    \includegraphics[height=3.5cm]{RR-icon}
  \end{wrapfigure}
  \vspace{0.5em}
  \begin{itemize}
  \setlength{\itemindent}{2em}
  {\Large
    \item[\ding{51}] 客観性
    \item[\ding{51}] 明瞭性
    \item[\ding{51}] 妥当性
  }\end{itemize}}
%#######*****#######*****#######*****#######***
\frame{
  \frametitle{研究の流れ}
  統計解析と論文の執筆は鬼門
  \begin{center}
  \includegraphics[width=11cm]{research-flow}
  \end{center}}

{\usebackgroundtemplate{\includegraphics[width=\paperwidth]{misc-files}}%
\frame{}} %解析フォルダは魑魅魍魎の世界

{\usebackgroundtemplate{\includegraphics[width=\paperwidth]{old-files-collection}}%
\frame{}} %過去のファイル

\frame{
  \frametitle{GitHubの特徴を活かしたRRな研究}
  ソーシャルコーディングの概念を研究に...
  \begin{center}
  \includegraphics[height=3.0cm]{research-flow-branch}
  \end{center}
  \begin{itemize}
    \item テキストファイルをもとにした進行
    \item 進捗状況に応じたブランチの切り替え
    \item 用途を記述したコミットコメント
    \item バックアップ、差分確認は気にしない
  \end{itemize}  
}

%\frame{
%  \frametitle{RRを行う際のガイドライン}
%  \vspace{-1.0em}
%  \begin{center}
%  \includegraphics[height=4cm]{course-and-guidelines}\\
%  {\Huge なければ作れば\\いいじゃない}\\
%  言い出しっぺの法則
%  \end{center}
%}

\frame{
  \frametitle{データ解析: \includegraphics[height=0.7cm]{r-logo}}
  \vspace{0.5em}
  リポジトリ内のソースコード、データファイルを実行すれば同様の結果が得られる
  \begin{wrapfigure}{o}[40pt]{60mm}
    \includegraphics[height=4.5cm]{r-code-markdown}
  \end{wrapfigure}
  \vspace{0.5em}
  \begin{itemize} 
    \item `Rstudio.app`で`.Rmd`, `.md`, `.html`ファイルを作成すれば、コードと作図の確認が可能
    \item 修正点はコミットコメントで
  \end{itemize}}

\frame{
  \frametitle{論文執筆: \includegraphics[height=0.7cm]{latex-logo}}
  テキストファイルなので差分の確認が容易。必要があればバックアップから復元
  \begin{wrapfigure}{o}[40pt]{60mm}
    \includegraphics[height=4.5cm]{latex-revise}
  \end{wrapfigure}
  \vspace{0.5em}
  \begin{itemize}
    \item 改訂を受けたらブランチを切り替え
    \item 「窓」からの脱出
  \end{itemize}  
}


\frame{
  \frametitle{Pull Request, Issuesを使った査読}
  \begin{center}
  \includegraphics[height=3cm]{instruction}\\
  \vspace{1.0em}
  \includegraphics[height=3.0cm]{research-flow-branch}
  \end{center}
}



\frame{
  \frametitle{今後の展望と課題}
  \begin{center}
  \includegraphics[height=5cm]{issues}
  \end{center}

}




\end{document}
